%-----------------------------------------------------------------------
\section{Dataset and Preprocessing}
%-----------------------------------------------------------------------

\subsection{Source and Scale}

Our experiments use ShuttleSet, a badminton broadcast dataset with stroke-level
annotations. Table~\ref{tab:data} summarizes the local snapshot: 44 raw
broadcast videos, 104 set annotation files, 36{,}482 annotated strokes, and
3{,}683 annotated rallies, along with six pre-extracted feature variants
comprising 33{,}481 prepared clips each. The active BST feature variant employs
a fixed sequence length of 100 frames. Videos~1--40 are reserved for
development; videos~41--44 form the held-out evaluation set used throughout
this work. While newer datasets such as FineBadminton provide richer semantic
annotations \cite{finebadminton}, ShuttleSet remains the only public resource
with aligned hit-frame indices, homography parameters, and stroke-level labels
required for our pipeline.

\begin{table}[t]
  \caption{Local ShuttleSet snapshot and data split used for evaluation.}
  \label{tab:data}
  \centering
  \small
  \begin{tabular}{lr}
    \toprule
    Item & Count or split \\
    \midrule
    Raw broadcast videos & 44 \\
    Set annotation files & 104 \\
    Annotated strokes    & 36{,}482 \\
    Annotated rallies    & 3{,}683 \\
    Prepared clips per variant & 33{,}481 \\
    Training videos      & 1--40 (train split) \\
    Validation videos    & 1--40 (val split, model selection only) \\
    Unseen evaluation    & 41--44 (held-out, never seen during training) \\
    Clean unseen predictions & 3{,}178 \\
    \bottomrule
  \end{tabular}
\end{table}

\subsection{Data Splits and Evaluation Protocol}
\label{sec:splits}

All 44 videos originate from ShuttleSet's pre-defined splits.
Videos~1--40 provide the \textbf{training set} (25{,}741 clips) and
\textbf{validation set} (4{,}241 clips); the validation set is used
\emph{only} for model selection---early stopping halted training at
epoch~6 based on validation macro-F1.
Videos~41--44 form the \textbf{held-out test set} (3{,}499 collated
clips; 3{,}178 clean primary events after removing feature-alias
duplicates). These four matches were withheld from all training, tuning,
and threshold decisions and evaluated exactly once to produce the figures
reported throughout this paper.

\subsection{Annotations and Label Preparation}

Table~\ref{tab:annotation_layers} summarizes the annotation layers used by the
pipeline. These annotations provide the temporal, spatial, and identity
information needed to convert raw video into ordered stroke-level examples.

\begin{table*}[h]
  \caption{Annotation layers used for feature construction, evaluation, and tactical mining.}
  \label{tab:annotation_layers}
  \centering
  \small
  \begin{tabular}{p{3.2cm} p{5.3cm} p{4.4cm}}
    \toprule
    Annotation layer & What it records & Role in the pipeline \\
    \midrule
    Stroke labels
      & Hitter side and stroke type for each event
      & Defines the supervised classification target \\[3pt]

    Rally identifiers
      & Rally ID and event order within each rally
      & Builds ordered stroke sequences for transitions and motifs \\[3pt]

    Hit-frame indices
      & Video frame where each stroke occurs
      & Anchors fixed-length stroke windows in raw video time \\[3pt]

    Homography parameters
      & Mapping from image coordinates to normalized court coordinates
      & Supports court normalization and movement profiling \\[3pt]

    Player identities
      & Persistent player information from ShuttleSet metadata
      & Used for evaluation and player-level movement analysis \\
    \bottomrule
  \end{tabular}
\end{table*}

The classification target contains 25 side-aware classes: one background
\texttt{none} class, 12 Top-side stroke classes, and 12 Bottom-side stroke
classes. Here, \emph{Top} and \emph{Bottom} are camera-relative sides, not
persistent player identities. For example, a Top-side smash means that the
far-court player, as seen by the broadcast camera, performed a smash.

\subsection{Preprocessing Pipeline}

Raw broadcasts are converted into aligned stroke windows through four sequential
stages, summarised in Table~\ref{tab:preproc}. Rally filtering identifies active
play segments and discards non-rally footage. Shuttle tracking predicts
frame-level coordinates, suppresses implausible jumps, and fills short gaps.
Window construction extracts a fixed-length frame range around each hit-frame,
and pose extraction obtains body keypoints and positions for both players. Each
stroke window includes four feature streams (player keypoints, positions,
shuttle coordinates, and validity masks) that feed the stroke classifier.

\begin{table*}[h]
  \caption{Preprocessing stages with corresponding operations and outputs.}
  \label{tab:preproc}
  \centering
  \small
  \begin{tabular}{p{3.3cm} p{5.5cm} p{4.1cm}}
    \toprule
    Stage & Operation & Output \\
    \midrule
    Rally filtering
      & Separate rally frames from non-play broadcast segments & Rally frame intervals \\
    Shuttle tracking
      & Predict shuttle coordinates, remove jumps, interpolate gaps & Denoised shuttle positions \\
    Window construction
      & Extract fixed-length frame range around each hit frame & Stroke windows \\
    Pose extraction
      & Extract, order, and normalize player keypoints & Player features and validity masks \\
    \bottomrule
  \end{tabular}
\end{table*}